\subsection{Quadratic forms.}

DEFINITION 2. --- Suppose the ring A is commutative. A map Q from E to A is said to be a quadratic form on E if

\begin{enumerate}
\def\labelenumi{\arabic{enumi})}
\item
  one has Q(αx) = α²Q(x) for α ∈ A and x ∈ E ;
\item
  the map \(\Phi : (x, y) \to \mathrm{Q}(x + y) - \mathrm{Q}(x) - \mathrm{Q}(y)\) of \(\mathbf{E} \times \mathbf{E}\) in A is a bilinear form.
\end{enumerate}

The bilinear form \(\Phi\) (which is necessarily symmetric) is called the bilinear form associated with Q. If \(\Phi\) is nondegenerate, Q is said to be nondegenerate.

Since Q(2x) = 4Q(x), it follows immediately from 2) that we have

\[
\Phi (x, x) = 2 \mathrm{Q} (x).\tag{13}
\]

In particular, if A is a ring of characteristic 2, the form \(\Phi\) is alternating.

We will say that two elements (resp. two subsets) of E are orthogonal with respect to Q if they are orthogonal with respect to the associated bilinear form Φ.

Let \((x_{i})_{i\in I}\) be a family of elements of E and \((a_{i})_{i\in I}\) a family of elements of A, all but finitely many of which are zero. By induction on the number of indices \(i\) for which \(a_{i}\neq0\) one shows easily that we have

\[
\mathrm{Q} (\sum_ {i} a _ {i} x _ {i}) = \sum_ {i} a _ {i} ^ {2} \mathrm{Q} (x _ {i}) + \sum_ {\{i, j \}} a _ {i} a _ {j} \Phi (x _ {i}, x _ {j}),\tag{14}
\]

the last summation being extended to the two-element subsets of I.

For every bilinear form \(f\) on \(\mathbf{E} \times \mathbf{E}\) one defines a quadratic form \(\mathbf{Q}\) by setting \(\mathrm{Q}(x) = f(x, x)\) ; the bilinear form \(\Phi\) associated with \(\mathbf{Q}\) is then defined by \(\Phi(x, y) = f(x, y) + f(y, x)\) for \(x, y\) in \(\mathbf{E}\) . Moreover, if one assumes that the scalar 2 has an inverse \(\frac{1}{2}\) in \(\mathbf{A}\) , there exists one and only one symmetric bilinear form \(f\) such that \(\mathrm{Q}(x) = f(x, x)\) , namely \(f = \frac{1}{2}\Phi\) ; the discriminant of \(f\) with respect to a system \(S = (x_1, \ldots, x_n)\) is also called the discriminant of \(\mathbf{Q}\) with respect to \(S\) . In this case, there is therefore a one-to-one correspondence between quadratic forms and symmetric bilinear forms on \(\mathbf{E}\) (cf. exercise 6).

In the case of a free module, we also have the following result:

PROPOSITION 2. --- Suppose that A is commutative and that E admits a basis \((e_i)_{i \in I}\) . Then, for every quadratic form Q on E, there exists a bilinear form \(f\) on E \(\times\) E such that Q(x) = f(x, x) for all \(x \in E\) . For every family \((b_{ij})_{(i,j) \in I \times I}\) of scalars such that \(b_{ij} = b_{ji}\) for \((i, j) \in I \times I\) , there exists one and only one quadratic form Q such that

\[
\mathrm{Q} (e _ {i}) = b _ {i i}, \quad \Phi (e _ {i}, e _ {j}) = b _ {i j} \text {   for   } i \neq j,\tag{15}
\]

where Φ denotes the bilinear form associated with Q; then Q is given by the formula

\[
\mathrm{Q} (\sum_ {i} a _ {i} e _ {i}) = \sum_ {\{i, j \}} b _ {i j} a _ {i} a _ {j},\tag{16}
\]

the last summation being extended over the subsets \(\{i, j\}\) of I having one or two elements.

Indeed, since formula (16) is only a transcription of formula (14), the uniqueness of a quadratic form Q satisfying (15) is proved. To prove its existence, we first note that there exists a family ( \(b_{ij}'\) ) of elements of A such that \(b_{ii}' = b_{ii}\) and \(b_{ij}' + b_{ji}' = b_{ij}\) for \(i \neq j\) ; for example, one obtains such a family by endowing I with a structure of totally ordered set (Ens., chap.~III, § 2, no. 3, th. 1) and setting \(b_{ij}' = b_{ij}\) for \(i<j\) and \(b_{ij}^{\prime}=0\) for \(i>j\) . Since the \(e_{i}\) form a basis of E, there exists a bilinear form \(f\) on \(\mathbf{E}\times\mathbf{E}\) such that \(f(e_{i},e_{j})=b_{ij}^{\prime}\) ; by setting \(\mathrm{Q}^{\prime}(x)=f(x,x)\) and denoting by \(\Phi^{\prime}\) the bilinear form associated with the quadratic form \(\mathrm{Q}^{\prime}\) , we obtain \(\mathrm{Q}^{\prime}(e_{i})=b_{ii}\) and \(\Phi^{\prime}(e_{i},e_{j})=f(e_{i},e_{j})+f(e_{j},e_{i})=b_{ij}\) . This proves our second assertion. As for the first, it follows immediately, because, by virtue of uniqueness, if a quadratic form Q satisfies (15), we have \(\mathrm{Q}(x)=\mathrm{Q}^{\prime}(x)=f(x,x)\) .

The module E equipped with the structure defined by a quadratic form Q is called a quadratic module. A homomorphism of the quadratic module (E, Q) into a quadratic module (E', Q') is a linear map u from E into E' such that Q = Q' ∘ u; if Φ and Φ' are the bilinear forms associated with Q and Q', we then have Φ(x, y) = Φ'(u(x), u(y)) for x ∈ E, y ∈ E; in other words, Φ' is the inverse image of Φ under u (§ 1, no. 1). Two quadratic forms Q and Q' on two A-modules E and E' are said to be equivalent if the corresponding quadratic modules are isomorphic.

Let \((\mathbf{E}_{i}, \mathbf{Q}_{i})_{i\in\mathbb{I}}\) be a family of quadratic modules, and let E be the direct sum of the modules \(E_{i}\) . One calls the external direct sum of the quadratic modules \((\mathbf{E}_{i}, \mathbf{Q}_{i})\) the quadratic module obtained by equipping E with the quadratic form Q defined by \(\mathrm{Q}(\sum_{i}x_{i}) = \sum_{i}\mathrm{Q}_{i}(x_{i})\) for \(x_{i} \in E_{i}\) . One also says that the quadratic form Q is the external direct sum of the quadratic forms \(Q_{i}\) .

If the forms \(Q_{i}\) are nondegenerate, so is Q.

Let Q be a quadratic form on the A-module E; if F is a submodule of E and if Q is constant on each class modulo F, the map \(\overline{Q}\) from E/F into A deduced from Q by passing to the quotient is evidently a quadratic form, and the canonical map from E onto E/F is a homomorphism for the structures of quadratic modules. For Q to be constant on each class modulo F, it is necessary and sufficient that one have \(Q(x + y) = Q(x)\) for \(x \in E\) and \(y \in F\) , that is, denoting by \(\Phi\) the bilinear form associated with Q, that one have \(Q(y) + \Phi(x, y) = 0\) for \(y \in F\) and \(x \in E\) . Setting x = 0, one sees that one has \(Q(y) = 0\) for \(y \in F\) , and hence \(\Phi(x,y)=0\) for \(x\in E\) and \(y\in F\) . In other words, if one calls the kernel of the quadratic module (E, Q) the set N of elements x of E such that \(Q(x)=0\) and \(\Phi(x,z)=0\) for all \(z\in E\) , for Q to be constant on each class modulo F, it is necessary and sufficient that F be contained in the kernel N of (E, Q). One verifies without difficulty that N is a submodule of E. For Q to be constant on each class modulo F, it is therefore necessary and sufficient that F be generated by elements of N.

One sees immediately that the kernel of the quadratic module E/N is \{0\}.

PROPOSITION 3. --- Let h be a homomorphism of A into a commutative ring A'. For every quadratic form Q on the A-module E, there exists one and only one quadratic form Q' on the A'-module A'⊗A E (chap.~III, 2nd ed., App. II, no. 10) such that one has

\[
\mathrm{Q} ^ {\prime} (1 \otimes x) = h (\mathrm{Q} (x))\tag{17}
\]

for all \(x \in \mathbf{E}\) . Moreover, the bilinear form \(\Phi'\) associated with Q' is obtained by extension of the ring of scalars from the bilinear form \(\Phi\) associated with Q.

We first show that, if there exists a quadratic form Q' satisfying (17), it is unique and the bilinear form \(\Phi'\) associated with Q' is obtained by extension of the ring of scalars from the form \(\Phi\) associated with Q. Indeed, this latter assertion follows from the fact that one has

\[
\begin{array}{r l} \Phi^ {\prime} (1 \otimes x, 1 \otimes y) & = Q ^ {\prime} (1 \otimes x + 1 \otimes y) - Q ^ {\prime} (1 \otimes x) - Q ^ {\prime} (1 \otimes y) \\ & = h (\Phi (x, y)) \end{array}\tag{18}
\]

for \(x \in \mathrm{E}, y \in \mathrm{E}\) . Formula (14) then shows that one has

\[
\mathrm{Q} ^ {\prime} \left(\sum_ {i} a _ {i} ^ {\prime} \otimes x _ {i}\right) = \sum_ {i} a _ {i} ^ {\prime 2} h (\mathrm{Q} (x _ {i})) + \sum_ {\{i, j \}} a _ {i} ^ {\prime} a _ {j} ^ {\prime} h (\Phi (x _ {i}, x _ {j}))\tag{19}
\]

for \(a_{i}^{\prime}\in\mathbf{A}^{\prime}\) and \(x_{i}\in\mathbf{E}\) , which proves the uniqueness of \(\mathbf{Q}^{\prime}\) .

To show the existence of Q', we shall first suppose that the module E admits a basis \((e_{i})_{i\in I}\) . There then exist elements \(b_{ij}\) of A such that \(b_{ij}=b_{ji}\) and \(\mathrm{Q}(\sum_{i}a_{i}e_{i})=\sum_{\{i,j\}}b_{ij}a_{i}a_{j}\) for \(a_{i}\in A\) (prop. 2). Since the elements \(1\otimes_{A}e_{i}\) form a basis of the

A'-module A' ⊗\_A E, one defines a quadratic form Q' on this latter module by setting

\[
\mathrm{Q} ^ {\prime} (\sum_ {i} a _ {i} ^ {\prime} \otimes e _ {i}) = \sum_ {\{i, j \}} a _ {i} ^ {\prime} a _ {j} ^ {\prime} h (b _ {i j})\tag{20}
\]

for \(a_i' \in A'\) ; whence, for \(x = \sum_{i} a_i e_i \in E\)

\[
Q ^ {\prime} (1 \otimes x) = Q ^ {\prime} (\sum_ {i} h (a _ {i}) \otimes e _ {i}) = \sum_ {\{i, j \}} h (a _ {i}) h (a _ {j}) h (b _ {i j}) = h (Q (x)),\tag{21}
\]

which proves the existence of Q' in this case.

Now let us turn to the general case. Let \((x_{i})_{i\in I}\) a system of generators of E, \(\mathbf{A}^{(1)}\) the module of formal linear combinations of elements of I (chap.~II, § 1, no. 8), and \((e_{i})_{i\in I}\) the canonical basis of \(\mathbf{A}^{(1)}\) . The linear map \(u\) of \(\mathbf{A}^{(1)}\) in E defined by \(u(e_{i}) = x_{i}\) is surjective since the elements \(x_{i}\) generate E. It follows (chap.~III, 2nd ed., App. II, no. 5, prop. 4) that the map \(1\otimes u\) of \(\mathbf{A}'\otimes\mathbf{A}^{(1)}\) in \(\mathbf{A}'\otimes\mathbf{E}\) is surjective, and its kernel P' is generated by elements of the form \(1\otimes p\) with \(u(p)=0\) Let then \(\mathrm{Q}_{1}^{\prime}\) be the extension to \(\mathbf{A}'\otimes_{\mathbf{A}}\mathbf{A}^{(1)}\) of the quadratic form \(\mathrm{Q}_{1}=\mathrm{Q}\circ u\) on \(\mathbf{A}^{(1)}\) . If \(p\) is an element of \(\mathbf{A}^{(1)}\) such that \(u(p)=0\) , one has \(\mathrm{Q}_{1}^{\prime}(1\otimes p)=h(\mathrm{Q}_{1}(p))=0\) and (denoting by \(\Phi_{1}^{\prime}\) the bilinear form associated with \(\mathrm{Q}_{1}^{\prime}\) ) \(\Phi_{1}^{\prime}(1\otimes p,1\otimes x)=h(\Phi(u(p),u(x)))=0\) for all \(x\in\mathbf{A}^{(1)}\) . Therefore, if \(u(p)=0\) ( \(p\in\mathbf{A}^{(1)}\) ), then \(1\otimes p\) belongs to the kernel of the quadratic module \(\mathbf{A}'\otimes_{\mathbf{A}}\mathbf{A}^{(1)}\) and consequently, as we saw above, there exists a quadratic form Q' on \(\mathbf{A}'\otimes_{\mathbf{A}}\mathbf{E}\) such that \(\mathrm{Q}_{1}^{\prime}=\mathrm{Q}^{\prime}\circ(1\otimes u)\) . Since \(u\) is surjective, we see that Q' satisfies condition (17). QED.

The quadratic form Q′, whose existence and uniqueness are guaranteed by Prop. 3, is called the extension of Q to A′ ⊗\_A E (with respect to h). One also says that Q' is obtained from Q by extension of the ring of scalars.

Exercises. --- 1) Let A be a field, E a left vector space over A, \(\Phi\) a sesquilinear form on E (for an antiautomorphism J of A). We assume that the rank (finite or infinite) of \(\Phi\) is \(\geqslant 2\) and that the relations \(\Phi(x, y) = 0\) and \(\Phi(y, x) = 0\) are equivalent.

\begin{enumerate}
\def\labelenumi{\alph{enumi})}
\item
  Show that there exists \(\lambda\neq0\) in A such that one has \(\Phi(y,x)=\lambda(\Phi(x,y))^{\mathfrak{J}}\) (Use exerc. 8 of § 1).
\item
  Show that there exists \(\alpha\in A\) such that the sesquilinear form \(\Phi\alpha\) (for the antiautomorphism \(\xi\to\alpha^{-1}\xi^{J}\alpha\) ) is Hermitian or alternating.
\end{enumerate}

(First note that one has \(\xi^{J^2} = \lambda^{-1}\xi\lambda^{-J}\) and \(\lambda\lambda^J = \lambda^J\lambda = 1\) . Then distinguish two cases according as \(\xi + \xi^J\lambda^{-1} = 0\) for all \(\xi \in A\) or not; in the second case, show that every element \(\neq 0\) of the form \(\alpha = \xi + \xi^J\lambda^{-1}\) answers the question).

\begin{enumerate}
\def\labelenumi{\arabic{enumi})}
\setcounter{enumi}{1}
\tightlist
\item
  Let \(\Phi\) be a sesquilinear form \(\varepsilon\) -hermitian on a finite-dimensional vector space E over a field A.
\end{enumerate}

\begin{enumerate}
\def\labelenumi{\alph{enumi})}
\item
  Show that for every vector subspace M of E, we have \((\mathbf{M}^{\mathbf{0}})^{\mathbf{0}} = \mathbf{M} + \mathbf{E}^{\mathbf{0}}\) and \(\dim M^{\mathbf{0}} + \dim M = \dim E + \dim (M \cap E^{\mathbf{0}})\) .
\item
  If \(M_{1}\) , \(M_{2}\) are two vector subspaces of E, show that one has dim \((\mathbf{M}_{1} \cap \mathbf{M}_{2}^{0}) + \dim (\mathbf{M}_{2} + \mathbf{M}_{1}^{0}) = \dim \mathbf{E} + \dim (\mathbf{M}_{1} \cap \mathbf{E}^{0})\) (consider the canonical map from E to \(E/E^{0}\) ).
\end{enumerate}

\begin{enumerate}
\def\labelenumi{\arabic{enumi})}
\setcounter{enumi}{2}
\tightlist
\item
  Let K be a commutative ring, \(f\) a monic polynomial in K{[}X{]}, of degree \(n \geqslant 1\) ; let A be the quotient algebra K{[}X{]}/(f) admitting as a basis over K the elements 1, \(\xi, \xi^2, \ldots, \xi^{n-1}\) (chap.~IV, § 1, no. 5, prop. 4). Show that the data of an A-module E is equivalent to the data of a K-module \(E_0\) and of a K-endomorphism \(j\) of \(E_0\) such that \(f(j) = 0\) . For every
\end{enumerate}

\(u \in \mathrm{E}^*\) , set \(u(x) = \sum_{k=0} u_k(x)\xi^k\) ; show that if \(\alpha_0 = f(0)\) is invertible in K, the map \(u \to u_0\) is a K-isomorphism of \(\mathbf{E}^*\) on \((\mathbf{E}_0)^*\) whose reciprocal isomorphism will be made explicit.

¶ 4) a) Let A be a ring (commutative or not), σ an automorphism of A such that there exists an invertible element γ ∈ A satisfying γ\^{}σ = γ, and such that one has ξ\^{}σ² = γξγ\^{}-1 for all ξ ∈ A. Let B be a left A-module having a basis of two elements (e₁, e₂); show that one defines on B a ring structure by taking as multiplication in B the composition law

\[
(\xi e _ {1} + \eta e _ {2}) (\xi^ {\prime} e _ {1} + \eta^ {\prime} e _ {2}) = (\xi \xi^ {\prime} + \eta \eta^ {\prime \sigma} \gamma) e _ {1} + (\eta \xi^ {\prime \sigma} + \xi \eta^ {\prime}) e _ {2}.
\]

For this ring structure, \(e_1\) is the identity element (identified with the identity element 1 of A); if we set \(e_2 = \rho\) , one has \(\rho^2 = \gamma\) and \(\rho\xi = \xi^\sigma\rho\) for all \(\xi \in A\) . Moreover, B is a right A-module, of which 1 and \(\rho\) form a basis. If A is a field, a necessary and sufficient condition for B to be a field is that \(\gamma\) is not of the form \(\lambda^\sigma\lambda\) (where \(\lambda \in A\) ). (Cf. Chap. VIII, § 12, Exerc. 8).

\begin{enumerate}
\def\labelenumi{\alph{enumi})}
\setcounter{enumi}{1}
\tightlist
\item
  Let J be an involutive antiautomorphism of A. Suppose there exists an invertible element \(\delta\in\mathsf{A}\) satisfying the following conditions:
\end{enumerate}

\[
\delta^ {J} = \delta , \quad \delta^ {\sigma} \delta = \gamma \gamma^ {J}, \quad (\xi^ {J}) ^ {\sigma} = \delta (\xi^ {\sigma}) ^ {J} \delta^ {- 1} \quad \text { for   all } \quad \xi \in A.\tag{1}
\]

Show that one can then extend J to an involutive antiautomorphism of B (still denoted J) by setting

\[
(\xi + \eta \rho) ^ {J} = \xi^ {J} + \gamma^ {- 1} \delta^ {\sigma} (\eta^ {J}) ^ {\sigma} \rho .\tag{2}
\]

If in addition A and B are fields, and if \(\sigma\) is not an inner automorphism, show that conditions (1) are necessary for the existence of an extension of J to an involutive antiautomorphism of B (by setting \(\rho^{J} = \alpha + \beta\rho\) , with \(\alpha, \beta\) in A, show that one necessarily has \(\alpha = 0\) by writing the condition \((\rho\xi)^{J} = \xi^{J}\rho^{J}\) for all \(\xi \in A\) ).

\begin{enumerate}
\def\labelenumi{\alph{enumi})}
\setcounter{enumi}{2}
\tightlist
\item
  We assume that conditions (1) are satisfied and that the involutive antiautomorphism J is extended to B by (2). Let F be a unitary B-module, E the unitary A-module obtained from F by restricting the ring of scalars to A; the map \(j: x \to \rho x\) is then a bijective semilinear map from E onto itself, relative to the automorphism \(\sigma\) of A and such that \(j^{2}(x) = \gamma x\) . Let \(\Phi\) be a Hermitian form on F (for J); for \(x \in E\) , \(y \in E\) , set \(\Phi(x, y) = \Phi_{1}(x, y) + \Phi_{2}(x, y)\rho\) , where \(\Phi_{1}(x, y) \in A\) , \(\Phi_{2}(x, y) \in A\) . Show that \(\Phi_{1}\) is a Hermitian form on E (with respect to J), such that
\end{enumerate}

\[
\Phi_ {1} (j (x), j (y)) = (\Phi_ {1} (x, y)) ^ {\sigma} \delta ;
\]

we have \(\Phi_{2}(x,y)=\Phi_{1}(x,j(y))\delta^{-1}\) , \(\Phi_{2}\) is a sesquilinear form on E for the antiautomorphism (in general non-involutive) \(\xi\to(\xi^{\mathrm{J}})^{\sigma}\) , such that

\[
\Phi_ {2} (j (x), j (y)) = (\Phi_ {2} (x, y)) ^ {\sigma} \delta^ {\sigma}
\]

and

\[
\Phi_ {2} (y, x) = \gamma^ {\mathrm{J}} \delta^ {- 1} ((\Phi_ {2} (x, y)) ^ {\mathrm{J}}) ^ {\sigma}.
\]

Converses. For \(\Phi\) to be nondegenerate, it is necessary and sufficient that \(\Phi_{1}\) or \(\Phi_{2}\) be nondegenerate. Special case where B is a quaternion algebra over a commutative ring K, corresponding to a pair \((\alpha, \beta)\) of elements of K, \(\beta\) being invertible, A the subalgebra K + Ku of B, and J and \(\sigma\) the map \(\xi \to \bar{\xi}\) (chap.~II, § 7, no. 8).

\begin{enumerate}
\def\labelenumi{\alph{enumi})}
\setcounter{enumi}{3}
\tightlist
\item
  Let u be an automorphism of the A-module E. For u to be an automorphism of the B-module F, leaving the form \(\Phi\) invariant, it is necessary and sufficient that u satisfy any two of the following three conditions:
\end{enumerate}

\(1^{\circ} u\) leaves \(\Phi_{1}\) invariant;

\(2^{o} u\) leaves \(\Phi_{2}\) invariant;

\(3^{o} u\) commutes with \(j\) .

\begin{enumerate}
\def\labelenumi{\arabic{enumi})}
\setcounter{enumi}{4}
\tightlist
\item
  Let A be a commutative ring, E an A-module, \((x_{i})_{i\in I}\) a system of generators of E; let R be the submodule of the module \(\mathrm{A}^{(\mathrm{I})}\) formed by the elements \((y_{i})_{i\in\mathrm{I}}\) such that \(\sum_{i}y_{i}x_{i}=0\) , and let \((a_{\lambda})_{\lambda\in\mathrm{L}}\) be a system of generators of R (with \(a_{\lambda}=(a_{\lambda i})_{i\in\mathrm{I}}\) ). Let \((b_{ij})\) be a family of elements of A \((i\in\mathrm{I},j\in\mathrm{I})\) . For there to exist a quadratic form Q such that \(\mathrm{Q}(x_{i})=b_{ii}\) and \(\Phi(x_{i},x_{j})=b_{ij}\) for \(i\neq j\) ( \(\Phi\) denoting the bilinear form associated with Q), it is necessary and sufficient that \(b_{ij}=b_{ji}\) for all i, j in I, and that, for all \(\lambda\in L\) and \(i\in I\) one has
\end{enumerate}

\[
\sum_ {\{i, j \}} b _ {i j} a _ {\lambda i} a _ {\lambda j} = \sum_ {j \neq i} b _ {i j} a _ {\lambda j} + 2 b _ {i i} a _ {\lambda i} = 0;\tag{3}
\]

we then have \(\mathrm{Q}(\sum_{i}a_{i}x_{i})=\sum_{\{i,j\}}b_{ij}a_{i}a_{j}\) . Deduce from this a new proof of prop. 3 of no. 4. (Note that the \(x_{i}^{\prime}=1\otimes x_{i}\) form a system of generators of \(A^{\prime}\otimes_{A}E\) , and that the \(A^{\prime}\) -module \(A^{\prime}\otimes_{A}E\) is isomorphic to \(A^{\prime(I)}/R^{\prime}\) , where \(A^{\prime(I)}\) is identified with \(A^{\prime}\otimes_{A}A^{(I)}\) and \(R^{\prime}\) is generated by the image of R under the canonical map from \(A^{(I)}\) into \(A^{\prime(I)}\) ).

\begin{enumerate}
\def\labelenumi{\arabic{enumi})}
\setcounter{enumi}{5}
\tightlist
\item
  Let A be a commutative ring of characteristic 2, E a free A-module, \(\mathcal{A}\) (resp. \(\mathcal{S},\mathcal{Q}\) ) the A-module of alternating (resp. symmetric bilinear, quadratic) forms on E. We have \(\mathcal{A}\subset\mathcal{S}\) ; one also defines a linear map \(\omega\) of \(\mathcal{S}\) in \(\mathcal{Q}\) , and a linear map \(\theta\) of \(\mathcal{Q}\) in \(\mathcal{A}\) as follows: for every bilinear form \(\Phi\in\mathcal{S}\) , \(\omega(\Phi)\) is the quadratic form \(x\to\Phi(x,x)\) , and for every quadratic form \(Q\in\mathcal{Q}\) , \(\theta(Q)\) is the bilinear form associated with Q, which is alternating. Show that \(\omega(0)=\mathcal{A}\) , \(\theta(\mathcal{Q})=\mathcal{A}\) and \(\bar{\theta}(0)=\omega(\mathcal{S})\) .
\end{enumerate}

¶ 7) Let A be a commutative ring, and E, F two A-modules. A map Q from E to F is said to be quadratic if it satisfies the following conditions: \(1^{\circ} \mathrm{Q}(\alpha x) = \alpha^{2} \mathrm{Q}(x)\) for \(\alpha \in \mathrm{A}\) , \(x \in \mathrm{E}\) ; \(2^{\circ}\) the map \((x, y) \to \mathrm{Q}(x + y) - \mathrm{Q}(x) - \mathrm{Q}(y)\) of \(\mathrm{E} \times \mathrm{E}\) into F is bilinear. If \(f\) is a linear map from an A-module \(\mathrm{E}_{1}\) into E, \(\mathrm{Q} \circ f\) is a quadratic map from \(\mathrm{E}_{1}\) into F.

\begin{enumerate}
\def\labelenumi{\alph{enumi})}
\tightlist
\item
  Let E be an A-module, \(\mathrm{A}^{(\mathrm{E})}\) the module of formal linear combinations of elements of E with coefficients in A (chap. II, § 1, no. 8), and for every \(x \in \mathrm{E}\) , let \(\varepsilon_x\) be the corresponding element of the standard basis of \(\mathrm{A}^{(\mathrm{E})}\) . Let \(\Gamma^2(\mathrm{E})\) be the quotient of \(\mathrm{A}^{(\mathrm{E})} \times (\mathrm{E} \otimes_{\mathbf{A}} \mathrm{E})\) by the submodule R generated by the elements ( \(\varepsilon_{x+y} - \varepsilon_x - \varepsilon_y, -x \otimes y\) ) and ( \(\varepsilon_{\lambda x} - \lambda^2\varepsilon_x, 0\) ), for \(x \in \mathrm{E}\) , \(y \in \mathrm{E}\) , \(\lambda \in \mathrm{A}\) . For every \(x \in \mathrm{E}\) , set \(\gamma(x) = \varphi(\varepsilon_x, 0)\) , denoting by \(\varphi\) the canonical mapping from \(\mathrm{A}^{(\mathrm{E})} \times (\mathrm{E} \otimes \mathrm{E})\) on \(\Gamma^2(\mathrm{E})\) ; we say that \(\gamma\) is the canonical mapping from E into \(\Gamma^2(\mathrm{E})\) . Show that \(\gamma\) is a quadratic map from E to \(\Gamma^2(\mathrm{E})\) and that, for every quadratic map Q from E into an A-module F, there exists one and only one linear map \(q\) of \(\Gamma^2(\mathrm{E})\) in F such that \(Q = q \circ \gamma\) (in other words, \((\Gamma^2(\mathrm{E}), \gamma)\) is a solution of a universal mapping problem; cf.~Sets, chap.~IV, § 3).
\end{enumerate}

For every pair of A-modules E, E′ and every linear map f from E into E′, show that, if γ′ denotes the canonical map from E′ into Γ²(E′), there exists one and only one linear map \(\bar{f}\) from Γ²(E) to Γ²(E') such that γ'∘f = \(\bar{f} \circ \gamma\) .

\begin{enumerate}
\def\labelenumi{\alph{enumi})}
\setcounter{enumi}{1}
\item
  Assume that E is the direct sum of two submodules M, N. Define a canonical isomorphism of \(\Gamma^{2}(\mathrm{E})\) onto the direct sum of modules \(\Gamma^{2}(\mathrm{M})\) , \(\Gamma^{2}(\mathrm{N})\) and \(\mathbf{M} \otimes \mathbf{N}\) (show that this direct sum is a solution of the same universal mapping problem as \(\Gamma^{2}(\mathrm{E})\) ).
\item
  Let F be a submodule of E, j the canonical injection of F into E. Define a canonical isomorphism of \(\Gamma^{2}(\mathrm{E}/\mathrm{F})\) onto
\end{enumerate}

\[
\Gamma^ {2} (\mathrm{E}) / (\bar {j} (\Gamma^ {2} (\mathrm{F})) + \psi (\mathrm{E} \times \mathrm{F})),
\]

where \(\psi(x, y) = \varphi(0, x \otimes j(y))\) for \(x \in \mathrm{E}, y \in \mathrm{F}\) . (Same method).

\begin{enumerate}
\def\labelenumi{\alph{enumi})}
\setcounter{enumi}{3}
\item
  Let A' be a commutative ring, and h a homomorphism from A into A'. Define a canonical isomorphism of \(\Gamma^{2}(\mathbf{A}'\otimes_{\mathbf{A}}\mathbf{E})\) onto \(\mathbf{A}'\otimes_{\mathbf{A}}\symbf{\Gamma}^{2}(\mathbf{E})\) (same method).
\item
  There exists one and only one linear map \(s\) of \(\Gamma^{2}(\mathrm{E})\) into \(\mathrm{E}\otimes\mathrm{E}\) such that \(s(\gamma(x))=x\otimes x\) for all \(x\in\mathrm{E}\) ; show that if \(\mathrm{E}\) is a free module, s is an isomorphism onto the submodule of symmetric tensors of order 2 on E.
\item
  Assume that A = Z and that E is a finite cyclic group of order n. Show that \(\Gamma^{2}(E)\) is a cyclic group of order n if n is odd, and of order 2n if n is even. (First note that if a is a generator of E, \(\gamma(a)\) is a generator of \(\Gamma^{2}(E)\) , and that \(\gamma(-ha) = \gamma(ha)\) for every integer h; deduce from this that if n is odd, \(n\gamma(a) = 0\) by taking \(h = (n - 1)/2\) ; similarly show that \(2n\gamma(a) = 0\) if n is even. Finally, prove that if n is odd (resp. even), there exists a quadratic map Q from E into a cyclic group of order n (resp. 2n) sending a onto a generator of this group.
\end{enumerate}

\begin{enumerate}
\def\labelenumi{\arabic{enumi})}
\setcounter{enumi}{7}
\tightlist
\item
  Let A be a commutative field, and let E, F be two vector spaces over A. Let g be a map from E to F, such that there exist three maps a, b, c from E × E to F, satisfying the identity
\end{enumerate}

\[
g (\lambda x + \mu y) = \lambda^ {2} a (x, y) + \lambda \mu b (x, y) + \mu^ {2} c (x, y)\tag{4}
\]

for all x, y in E, λ, μ in A.

\begin{enumerate}
\def\labelenumi{\alph{enumi})}
\tightlist
\item
  Show that we have \(a(x, y) = g(x)\) , \(c(x, y) = g(y)\) , \(b(y, x) = b(x, y)\) and \(b(\lambda x, y) = \lambda b(x, y)\) ; moreover, if we set
\end{enumerate}

\[
d (x, y, z) = b (x + y, z) - b (x, z) - b (y, z)
\]

show that we have \(d(x,y,z)=d(y,z,x)=d(z,x,y)\) and conclude from this that \(d(\lambda x,\mu y,\nu z)=\lambda\mu\nu d(x,y,z)\) .

\begin{enumerate}
\def\labelenumi{\alph{enumi})}
\setcounter{enumi}{1}
\tightlist
\item
  From a), deduce that if \(\mathbf{A} \neq \mathbf{F}_2\) one necessarily has \(d(x, y, z) = 0\) and consequently that \(g\) is a quadratic map (exerc. 7). On the contrary, if \(\mathbf{A} = \mathbf{F}_2\) and if dim \(\mathrm{E} \geqslant 3\) show that there exist maps \(g\) from E to A, satisfying (4) and for which \(d(x, y, z)\) is not identically zero.
\end{enumerate}

\begin{enumerate}
\def\labelenumi{\arabic{enumi})}
\setcounter{enumi}{8}
\tightlist
\item
  Let A be a commutative ring, E an A-module having a basis of n elements, Φ a symmetric bilinear form on E. Let e be an element of \(\bigwedge^{n}\) E forming a basis of this module, \(\Delta\) the discriminant of \(\Phi\) with respect to \(e\) , \(\varphi_{p}\) the canonical isomorphism of \(\bigwedge^{p}\) E onto \(\bigwedge^{n-p}\) E* relative to \(e\) (chap.~III, §8, no. 5). For \(x \in \bigwedge^{p}\) E, let \(d_{(p)}(x) = \bar{\varphi}_{n-p}^{-1}(d_{\Phi(p)}(x))\) ; show that the map \(d_{(p)}\) of \(\bigwedge^{p}\) E into \(\bigwedge^{n-p}\) E possesses the following properties:
\end{enumerate}

\begin{enumerate}
\def\labelenumi{\alph{enumi})}
\item
  For every \(x \in \wedge \mathrm{E}\) , \(d_{(n - p)}(d_{(p)}(x)) = (-1)^{p(n - p)}\Delta x\) .
\item
  For every pair of elements \(x, y\) of \(\wedge\) E, we have
\end{enumerate}

\[
x \wedge d _ {(p)} (y) = \Phi_ {(p)} (x, y) e \quad \text { and } \quad \Phi_ {(n - p)} (d _ {(p)} (x), d _ {(p)} (y)) = \Delta \Phi_ {(p)} (x, y).
\]

\begin{enumerate}
\def\labelenumi{\alph{enumi})}
\setcounter{enumi}{2}
\item
  We further assume that A is a field and that \(\Phi\) is nondegenerate. Then, if \(x\) is a decomposable \(p\) -vector \(\neq 0\) corresponding to a subspace F of E (chap. III, § 7, no. 3), \(d_{(p)}(x)\) is a decomposable \((n-p)\) -vector \(\neq 0\) corresponding to the subspace F \(^{0}\) orthogonal to F.
\item
  Extend the previous results to the case where (A being a commutative ring), \(\Phi\) is a sesquilinear form \(\varepsilon\) -hermitian for an involutive automorphism \(J \neq 1\) of A.
\end{enumerate}

\begin{enumerate}
\def\labelenumi{\arabic{enumi})}
\setcounter{enumi}{9}
\tightlist
\item
  \begin{enumerate}
  \def\labelenumii{\alph{enumii})}
  \tightlist
  \item
    Let A be a commutative ring, E an A-module having a basis of 3 elements, \(\Phi\) a symmetric bilinear form on E. With the notations of exerc. 9, for any two elements \(x\) , \(y\) of E, one sets \(x \overline{\wedge} y = d_{(2)}(x \wedge y)\) , and one says that this element is the vector product
  \end{enumerate}
\end{enumerate}

\[
\stackrel {{3}} {\wedge} \mathrm{E})
\]

of \(x\) and of \(y\) (relative to \(\Phi\) and to the basis \(e\) of \(\wedge E\) ). Show that \((x, y) \to x \overline{\wedge} y\) is an alternating bilinear map from \(E \times E\) into \(E\) , and that \(x \overline{\wedge} y\) is orthogonal to \(x\) and to \(y\) .

\begin{enumerate}
\def\labelenumi{\alph{enumi})}
\setcounter{enumi}{1}
\tightlist
\item
  Let \(\alpha, \beta\) be two invertible elements of A, B the quaternion algebra over A corresponding to the pair \((\alpha, \beta)\) (chap.~II, § 7, no. 8), 1, u, v, w the canonical basis of B over A; let E be the submodule of B having u, v, w as basis. Show that if x, y are two quaternions belonging to E, one has
\end{enumerate}

\[
x y = \Phi (x, y) + x \overline {{{{\wedge}}}} y
\]

where \(\Phi\) is a symmetric bilinear form on E, such that the linear maps associated with \(\Phi\) are bijective, and \(x\overline{\wedge}y\) is the vector product of x and y relative to the form \(\Phi\) and to the basis \(\alpha^{-1}\beta^{-1}u\wedge v\wedge w\) of \(\bigwedge^{3}\mathbf{E}\) .

\begin{enumerate}
\def\labelenumi{\arabic{enumi})}
\setcounter{enumi}{10}
\tightlist
\item
  Let \(\Phi\) be a nondegenerate \(\varepsilon\) -hermitian sesquilinear form on a finite-dimensional vector space E. One says that a vector subspace M of E is weakly orthogonal to a vector subspace N (relative to \(\Phi\) ) if one of the two subspaces M, N° contains the other.
\end{enumerate}

\begin{enumerate}
\def\labelenumi{\alph{enumi})}
\item
  Show that the relation “M is weakly orthogonal to N” is symmetric.
\item
  If M and N are weakly orthogonal, \(M^{0}\) and \(N^{0}\) are weakly orthogonal.
\item
  If M and N are weakly orthogonal and if M ∩ N = \{0\}, M and N are orthogonal.
\end{enumerate}

